\section{Introduction}
A pseudo-Riemannian manifold is a smooth manifold $M$ equipped with
a smooth non-degenerate symmetric bilinear tensor $g$
 of signature $(p,q)$ on $M$.
It is called Riemannian if $q=0$, and Lorentzian if $q=1$.
As in the Riemannian case,
the Laplacian $\square_{M}:=\operatorname{div}_{M} \circ \operatorname{grad}_{M}$
is defined as a second-order differential operator on $M$. We~note that it is a
hyperbolic differential operator if $M$ is Lorentzian. We~write $L^2(M)$ for the Hilbert space of square-integrable functions on $M$
with respect to the Radon measure induced by the pseudo-Riemannian structure.
For $\lambda \in \mathbb{C}$, we denote by
\begin{gather*}
L^2_{\lambda}(M) :=
\big\{f \in L^2(M) \mid \square_{M}f = \lambda f \text{ in the weak sense}\big\}.
\end{gather*}
The set of $L^2$-eigenvalues
$\mathrm{Spec}_{d}(\square_{M}) := \big\{ \lambda \in \mathbb{C}
\mid L^2_{\lambda}(M)\neq0\big\}$
is called the \textit{discrete spectrum} of~$\square_{M}$.

Our interest is the multiplicities of $L^2$-eigenvalues $\lambda$ of $\square_{M}$,
denoted by
\begin{gather*}
\mathcal{N}_{M}(\lambda):= \operatorname{dim}_{\mathbb{C}}L^2_{\lambda}(M)
\in \mathbb{N}\cup\{\infty\}.
\end{gather*}
In the Riemannian case, the Laplacian is an elliptic differential operator
and the distribution of~its discrete spectrum has been investigated extensively,
such as the Weyl law for compact Riemannian manifolds.
However, it is not the case for non-Riemannian manifolds.
Kobayashi~\cite{Kob16}, and later Fox--Strichartz~\cite{FoSt18},
investigated the distribution of the discrete spectrum of
the Laplacian $\square_{M}$ of some pseudo-Riemannian manifolds, i.e., when
$M$ is the flat pseudo-Riemannian manifold $\mathbb{R}^{p,q}/\mathbb{Z}^{p+q}$
and is the Lorentzian manifold $S^1\times S^{q}$, respectively.

Let us recall some basic notions.
A \textit{discontinuous group} for a homogeneous manifold \mbox{$X\!=\!G/H$}
is a discrete subgroup $\Gamma$ of $G$ acting properly discontinuously
and freely on $X$ (Kobayashi \cite[Definition~1.3]{Ko01}). In~this case, the quotient space $X_{\Gamma}:=\Gamma\backslash X$ carries a
$C^{\infty}$-manifold structure such that
the quotient map $p_{\Gamma}\colon X\rightarrow X_{\Gamma}$ is a covering of
$C^{\infty}$ class, hence $X_{\Gamma}$ has a $(G,X)$-structure induced by $p_{\Gamma}$.
If we drop the assumption of freeness, $X_{\Gamma}$ is not always a manifold but
carries a nice structure called an orbifold or $V$-manifold.
Proper discontinuity is a more
serious assumption which assures $X_{\Gamma}$ to be Hausdorff in the quotient topology. We~remark that the action of a discrete subgroup $\Gamma$ on $X$
may fail to be properly discontinuous
when $H$ is noncompact. In~order to overcome this difficulty,
Kobayashi~\cite{Kob96} and Benoist~\cite{Ben96} established the properness criterion
for reductive $G$ generalizing the original criterion by Kobayashi~\cite{Kob89}.
Whereas discontinuous groups for the de Sitter space
$\mathrm{dS}^{n}:=\mathrm{SO}_{0}(n,1)/\mathrm{SO}_{0}(n-1,1)$
are always finite groups
(the Calabi--Markus phenomenon, see~\cite{CaMa62, Kob89}),
there are a rich family of discontinuous groups for
the anti-de Sitter space, see, e.g.,~\cite{Go85, Kob98, Sa00}. We~treat, in this article, the three-dimensional anti-de Sitter space
$\mathrm{AdS}^{3}:=\mathrm{SO}_{0}(2,2)/(\{\pm1\}\times\mathrm{SO}_{0}(2,1))$.

For $m\in \mathbb{N}$, we set
\begin{gather*}%\label{lambda}
\lambda_{m} := 4m(m-1).
\end{gather*}
